% generated by roam.py handbook export -- не редактировать, править заметки
% 20260916214817-векторное_пространство.org (body 1522a4bcd559)

\begin{definition}[{Векторное пространство}]\label[definition]{note:1B626AA4-35C8-4C7C-9D3D-D8C26D98F40C}
  Пусть $K$~--- поле (в курсе $K = \R$, реже $\C$). Множество $V$ с операциями сложения $V \times V \rightarrow V$ и умножения на скаляр $K \times V \rightarrow V$ называется \textbf{векторным} (линейным) пространством над $K$, если $\forall x, y, z \in V\ \forall \alpha, \beta \in K$:

  \begin{enumerate}
    \item $(x + y) + z = x + (y + z)$
    \item $x + y = y + x$
    \item $\exists\, 0 \in V:\ x + 0 = x$
    \item $\forall x\ \exists\, (-x) \in V:\ x + (-x) = 0$
    \item $\alpha (\beta x) = (\alpha \beta) x$
    \item $1 \cdot x = x$
    \item $\alpha (x + y) = \alpha x + \alpha y$
    \item $(\alpha + \beta) x = \alpha x + \beta x$
  \end{enumerate}

  Примеры: $\R^n$; матрицы $\R^{m \times n}$; непрерывные функции $C[a;b]$; многочлены степени не выше $k$.
\end{definition}

\begin{remark}
  В курсе почти всегда $V = \R^n$. Для \hyperref[note:316B5BFA-8B9F-4267-9086-D6593AAB5F9B]{выпуклых множеств} важно, что поле вещественное: определение использует $\lambda \in [0;1]$ и комбинации $\lambda x + (1 - \lambda) y$.
\end{remark}

\begin{remark}
  Дополнительные структуры на $V$: \hyperref[note:F2942E98-6115-4600-AA3F-E0F687D9C732]{скалярное произведение} (и порожденная им норма $\norm{x} = \sqrt{\left\langle x; x \right\rangle}$), полнота по этой норме~--- \hyperref[note:5D688904-8445-47B5-9D39-8BCC39F9C354]{гильбертово пространство}.
\end{remark}
